%%%%%%%%%%%%%%%%%%%%%%%%%%%%%%%%%%%%%%%%%%%%%%%%
% COPYRIGHT: (C) 2012-2026 FAU FabLab and others
% CC-BY-SA 3.0
%%%%%%%%%%%%%%%%%%%%%%%%%%%%%%%%%%%%%%%%%%%%%%%%
%
% Die Hinweisschilder, eingebunden von hinweisschilder-a6.tex und
% hinweisschilder-a7.tex. Aufbau und Optionen: siehe schilder-layout.tex.
%
% \begin{schild}[farbe=..., modell=..., einweisung=REPO, kontakt=...]{Titel}{Status}
%   Text
% \end{schild}
%
% Farben (Hintergrund):
%   gruen      Benutzung mit allgemeiner Werkstatt-Einweisung
%   gelbgruen  wie gruen, aber empfindlich: Hinweise beachten
%   gelb       nur mit Einweisung für dieses Gerät
%   gelbrot    nur mit Einweisung, Aufsicht erforderlich
%   rot        nur nach Rücksprache bzw. gar nicht
%
% Gelbe Schilder mit leerem Status bekommen einheitlich "Benutzung nur mit
% unterschriebener Einweisung".
%
% QR-Code: einweisung=REPO zeigt auf das neueste Release der Einweisung
% (github.com/fau-fablab/REPO), sonst seite=... für fablab.fau.de.
%
% Im PDF werden die Schilder sortiert: erst Geräte mit eigener Einweisung,
% dann allgemeine Werkstatt-Einweisung, dann ohne Einweisung (je ab neuem Bogen).

\hyphenation{Be-treu-er Be-treu-er-in-nen Werk-statt-Ein-wei-sung Multi-funk-tions-tisch}

\newcommand{\werkstattEinweisung}{Benutzung mit allgemeiner Werkstatt-Einweisung}

% ---------------------------------------------------------------------
% Zerspanung
% ---------------------------------------------------------------------
\begin{schild}[farbe=gelbrot, einweisung=fraese-einweisung,
	kontakt={AG Zerspanung, \mail{zerspanung@fablab.fau.de}}]
	{CNC-Fräse}{Benutzung nur mit Einweisung, Aufsicht erforderlich!}
Bei Interesse frage die Betreuer:innen!

Je nach Einweisungsstufe sind bestimmte Tätigkeiten nur unter Aufsicht erlaubt.
\end{schild}

\begin{schild}[farbe=gelbrot, einweisung=drehbank-einweisung,
	kontakt={AG Zerspanung, \mail{zerspanung@fablab.fau.de}}]
	{CNC-Drehbank}{Benutzung nur mit Einweisung, Aufsicht erforderlich!}
Bei Interesse frage die Betreuer:innen!

Je nach Einweisungsstufe sind bestimmte Tätigkeiten nur unter Aufsicht erlaubt.
\end{schild}

% ---------------------------------------------------------------------
% Lasercutter, 3D-Drucker, Schneideplotter
% ---------------------------------------------------------------------
\begin{schild}[farbe=gelb, modell=LTT iLaser, einweisung=lasercutter-einweisung,
	kontakt={AG Technik, \mail{kontakt@fablab.fau.de}}]
	{Lasercutter}{}
Die Einweisung gilt nur für dieses Gerät (eigene Einweisungsliste).
Bei Interesse frage einfach die Betreuer:innen!

Einweisungen vor Oktober 2018 oder für den kleineren Lasercutter (Epilog Zing)
sind \textbf{hier nicht gültig}.
\end{schild}

\begin{schild}[farbe=gelb, modell=Epilog Zing, einweisung=lasercutter-einweisung,
	kontakt={AG Technik, \mail{kontakt@fablab.fau.de}}]
	{Lasercutter}{}
Die Einweisung gilt nur für dieses Gerät (eigene Einweisungsliste).
Bei Interesse frage einfach die Betreuer:innen!

Einweisungen für den großen Lasercutter (LTT iLaser) sind \textbf{hier nicht gültig}.
\end{schild}

\begin{schild}[farbe=gruen, seite=lasercutter]
	{Lasermaterial}{Selbstbedienung. Bitte zuerst die Reste aufbrauchen!}
Kleinere Reste befinden sich links in den Schubladen.
Angefangene Platten liegen oben auf den neuen.
Wenn kein Rest passt, fange eine neue Platte an.

\textbf{Die Preisliste hängt links neben der Tür.}
\end{schild}

\begin{schild}[farbe=gelb, modell=Bambu Lab, einweisung=3d-drucker-einweisung]
	{FDM-3D-Drucker}{}
Bei Interesse frage die Betreuer:innen!
\end{schild}

\begin{schild}[farbe=gelb, einweisung=sla-drucker-einweisung, symbole={M004, M009}]
	{SLA-Drucker}{}
Bei Interesse frage einfach die Betreuer:innen!
\end{schild}

\begin{schild}[farbe=gelb, modell=Formlabs Form Wash und Form Cure, einweisung=sla-drucker-einweisung,
	symbole={GHS05, GHS07, GHS08, GHS09, M004, M009}]
	{Waschen und Aushärten}{}
Es gilt die SLA-Einweisung.

Bei Interesse frage einfach die Betreuer:innen!
\end{schild}

\begin{schild}[farbe=gelb, einweisung=schneideplotter-einweisung]
	{Schneideplotter}{}
Bei Interesse frage einfach die Betreuer:innen!
\end{schild}

\begin{schild}[farbe=gelb, einweisung=schneideplotter-einweisung, symbole={W017}]
	{T-Shirt-Presse}{}
Es gilt die Schneideplotter-Einweisung.

Bei Interesse frage einfach die Betreuer:innen!
\end{schild}

\begin{schild}[farbe=gruen]
	{Plotterfolien}{Bitte genau rechtwinklig abschneiden!}
Reste sind in der Kiste über dem Plotter.
Verwende zum Abschneiden bitte die Schneidekante am Plotter.

Preis und Folientyp (Aufkleber/Textil-Flex/Flock) stehen in den Folienrollen auf einem Etikett.

\textbf{Fehlt was?} Fast ausverkauft? Schreib es auf die Nachkaufliste am Whiteboard!
Teilesuche online: \web{app.fablab.fau.de}
\end{schild}

% ---------------------------------------------------------------------
% Elektronik
% ---------------------------------------------------------------------
\begin{schild}[farbe=gruen, einweisung=werkstatt-einweisung]
	{Elektro-Arbeitsplätze}{Benutzung der Geräte mit allgemeiner Werkstatt-Einweisung}
Bei Interesse frage einfach die Betreuer:innen!

Arbeitsplatz \textbf{sauber und leer} verlassen: Werkzeug verräumen, Müll entsorgen,
Kabel zurücklegen. Liegengebliebene Sachen können einfach so verschwinden!
\end{schild}

\begin{schild}[farbe=rot, kontakt={Max B, Philipp, \mail{kontakt@fablab.fau.de}}]
	{Entlötstation}{Bitte frage die Betreuer:innen!}
Vorsicht, empfindlich: Nicht klopfen, nicht hebeln!

Spitzenwechsel, Reinigung, Entleeren usw. nur durch Gerätebetreuer:innen!
Bitte lasse dir die Benutzung von den Betreuer:innen erläutern.

Bei Problemen bitte an Max B, Philipp oder Patrick wenden.
\end{schild}

\begin{schild}[farbe=gelb, einweisung=platinenaetzer-einweisung,
	kontakt={Max, \mail{technik@fablab.fau.de}}, symbole={W023, M004, M009}]
	{Platinenfertigung}{}
Bei Interesse frage einfach die Betreuer:innen!
\end{schild}

\begin{schild}[farbe=gelb, einweisung=platinenaetzer-einweisung,
	kontakt={Max, \mail{technik@fablab.fau.de}}]
	{DuKo-Presse für Platinen}{}
Vorsichtig verwenden! Das Werkzeug ist sehr empfindlich!

Bei Interesse frage einfach die Betreuer:innen!
\end{schild}

\begin{schild}[farbe=rot, kontakt={Philipp, \mail{kontakt@fablab.fau.de}}]
	{Sprühätzer}{Darf nicht benutzt werden}
Es wurde noch keine Einweisung erstellt und freigegeben.
Bei Interesse wende dich bitte an Philipp!

\textbf{Bei Zuwiderhandlung erfolgt Lab-Verbot!}
\end{schild}

\begin{schild}[farbe=gelb, einweisung=reflow-ofen-einweisung,
	kontakt={Julian, \mail{julian.neureuther@fablab.fau.de}}]
	{Lötpastendrucker}{}
Es gilt die Reflow-Ofen-Einweisung.

Bei Interesse frage einfach die Betreuer:innen!
\end{schild}

\begin{schild}[farbe=gelb, modell=VaporPhase 1, einweisung=dampfphasenloetanlage-einweisung,
	symbole={W017, M004, P022}]
	{Dampfphasen\-lötanlage}{}
Bei Interesse frage die Betreuer:innen!
\end{schild}

\begin{schild}[farbe=rot, modell=Siglent SSA3021X, kontakt={Markus F+W, \mail{kontakt@fablab.fau.de}}]
	{Spektrumanalysator}{Benutzung nur nach Rücksprache}
\textbf{Achtung: ESD-gefährdet!}

Im Moment gibt es leider noch keine Einweisung zu diesem Gerät.
Bitte sprich mit den Betreuer:innen.
\end{schild}

\begin{schild}[farbe=gruen, einweisung=werkstatt-einweisung]
	{Tastköpfe}{\werkstattEinweisung}
Bei Unklarheiten frage die Betreuer:innen!

Kabel nicht knicken, nicht am Kabel ziehen.
Nach Benutzung vollständig zurücklegen.
\end{schild}

\begin{schild}[farbe=gelbgruen, modell={Testec, 250\,MHz / 20\,MHz, 1:10 / 1:1}, einweisung=werkstatt-einweisung]
	{Tastköpfe}{\werkstattEinweisung, bitte sorgfältig behandeln}
6 Stück mit Masseklemme und Hakenspitze, dazu Zubehör: Krokoklemme auf BNC,
BNC-BNC-Kabel, T-Stücke.

Teilerverhältnis (1:1 oder 1:10) auch am Oszilloskop einstellen, bei 1:10 den
Tastkopf vorher abgleichen. Nach Benutzung vollständig zurücklegen.
\end{schild}

\begin{schild}[farbe=gelbgruen, modell=bis 350\,MHz, einweisung=werkstatt-einweisung]
	{Hochwertige Tastköpfe}{\werkstattEinweisung, bitte Hinweise beachten}
Nur verwenden, wenn die Messung sie wirklich braucht, sonst die einfachen Tastköpfe nehmen.
\begin{itemize}
	\item Spitzen sind empfindlich: nicht verbiegen, nicht als Werkzeug benutzen
	\item Kabel nicht knicken, nicht am Kabel ziehen
	\item maximale Eingangsspannung beachten
	\item Masseklemme nur an Masse anschließen
	\item vor der Messung am Oszilloskop abgleichen
\end{itemize}
\end{schild}

% ---------------------------------------------------------------------
% Werkstatt
% ---------------------------------------------------------------------
\begin{schild}[farbe=gruen, einweisung=werkstatt-einweisung]
	{Werkbank}{Benutzung der Werkzeuge mit allgemeiner Werkstatt-Einweisung}
Bei Interesse frage einfach die Betreuer:innen!
Beachte auch die Aufschrift an der Standbohrmaschine!

Arbeitsplatz \textbf{sauber und leer} verlassen: Werkzeug verräumen, Müll entsorgen.
Liegengebliebene Sachen können einfach so verschwinden!
\end{schild}

\begin{schild}[farbe=rot, einweisung=werkstatt-einweisung]
	{Standbohrmaschine}{Benutzung nur nach Rücksprache}
Vor \textbf{jeder} Benutzung die Betreuer:innen fragen! Sie weisen dich ein, wenn du
noch keine Einweisung hast, prüfen dein Werkstück und geben es frei.
\begin{itemize}
	\item Haare und alles andere Lose gesichert?
	\item Werkstück festgespannt (nicht mit der Hand halten!)?
	\item Bohrer richtig eingespannt?
	\item Im Zweifelsfall: nicht freigeben!
\end{itemize}
\textbf{Bitte gut aufpassen!}
\end{schild}

\begin{schild}[farbe=gruen]
	{Druckluft}{Ventile stets vor dem Trennen einer Verbindung schließen!}
Verwende die Druckluftanlage zum Ausblasen oder Schlauch aufpumpen.

Beim Ausblasen von Spänen ggf. Schutzbrille tragen und nicht auf Mitmenschen zielen!
Bitte \textbf{nicht} in der Elektrowerkstatt ausblasen, Dreck sollte draußen bleiben!

Bei Problemen frage die Betreuer:innen!
\end{schild}

\begin{schild}[farbe=gruen, modell=Festool CLEANTEC CTM 26]
	{Absaugmobil}{Bitte gerne viel benutzen!}
Zum Nass- und Trockensaugen geeignet.
\end{schild}

\begin{schild}[farbe=rot, modell=Festool MFT 3, seite=multifunktionstisch,
	kontakt={FabLab-Aktive, \mail{technik@fablab.fau.de}}]
	{Multifunktionstisch}{Benutzung nur nach Rücksprache}
Im Moment gibt es leider noch keine Einweisung zu diesem Gerät.
Bitte sprich mit den Betreuer:innen.
\end{schild}

\begin{schild}[farbe=gelb, modell=Festool OF 1010, einweisung=festool-oberfraese-einweisung]
	{Oberfräse}{}
Bei Interesse frage die Betreuer:innen!
\end{schild}

\begin{schild}[farbe=gelb, modell=Festool DOMINO DF 500 RQ-Plus, einweisung=festool-duebelfraese-einweisung]
	{Dübelfräse}{}
Bei Interesse frage die Betreuer:innen!
\end{schild}

\begin{schild}[farbe=gelb, einweisung=shaper-origin-einweisung]
	{Shaper Origin}{}
Bei Interesse frage die Betreuer:innen!
\end{schild}

\newcommand{\erprobungsbetriebEins}{Verwendung nur durch stark eingeschränkten Nutzerkreis}
\newcommand{\erprobungsbetriebZwei}{Diese Maschine darf nur nach Unterschrift auf der
zugehörigen Liste verwendet werden. Einweisung unfertig, daher dürfen weitere Personen
nur nach ausführlichem Beschäftigen mit der Maschine und ihrer Anleitung und durch
Unterschrift von Philipp (\mail{philipp@fablab.fau.de}) auf die Liste hinzugefügt
werden. Sobald es eine Einweisung gibt, wird dieses Schild ersetzt und eine normale
Benutzung wird mit unterschriebener Einweisung möglich sein. Bei der Benutzung sind die
Sicherheitsregeln aus der Anleitung zu befolgen.}

\begin{schild}[farbe=rot, modell=Festool TS 55 R, einweisung=festool-kreissaege-einweisung,
	kontakt={Philipp, \mail{technik@fablab.fau.de}}]
	{Tauchsäge}{\erprobungsbetriebEins}
\erprobungsbetriebZwei
\end{schild}

\begin{schild}[farbe=rot, modell=Proxxon, kontakt={Philipp, \mail{kontakt@fablab.fau.de}}]
	{Tellerschleifgerät}{\erprobungsbetriebEins}
\erprobungsbetriebZwei
\end{schild}

\begin{schild}[farbe=rot, modell=Proxxon, kontakt={Philipp, \mail{kontakt@fablab.fau.de}}]
	{Schleif- und Poliergerät}{\erprobungsbetriebEins}
\erprobungsbetriebZwei
\end{schild}

\begin{schild}[farbe=rot, modell=Proxxon, einweisung=tischkreissaege-einweisung,
	kontakt={David, \mail{kontakt@fablab.fau.de}}]
	{Tischkreissäge}{Benutzung nur mit unterschriebener Einweisung, immer unter Aufsicht durch Betreuer:innen!}
Bei Interesse frage die Betreuer:innen!
\end{schild}

\begin{schild}[farbe=rot, modell=Festool, kontakt={David, \mail{kontakt@fablab.fau.de}}]
	{Tischkreissäge}{Benutzung nur mit unterschriebener Einweisung, immer unter Aufsicht durch Betreuer:innen!}
Bei Interesse frage die Betreuer:innen!
\end{schild}

% Festool-Handmaschinen und Zubehör (Issue #1)
% \festool[Symbole]{Gerät}{Modell}{Hinweise}
\newcommand{\festool}[4][]{%
\begin{schild}[farbe=gruen, modell={#3}, einweisung=werkstatt-einweisung, symbole={#1}]
	{#2}{\werkstattEinweisung}
Bei Unklarheiten frage die Betreuer:innen!

#4

Nach Benutzung vollständig in den Systainer zurücklegen.
\end{schild}}

\festool{Akkuschrauber}{Festool T 18+3}
	{Akkus nach Benutzung wieder aufladen.}

\festool[M004]{Stichsäge}{Festool CARVEX PS 420}
	{Werkstück fest einspannen.
	Passendes Sägeblatt für das Material verwenden.}

\festool{Zubehör Stichsäge}{Festool CARVEX PS 420}
	{Sägeblätter und Zubehör für die CARVEX.
	Stumpfe oder abgebrochene Sägeblätter bitte den Betreuer:innen melden.}

\festool{Schwingschleifer}{Festool RUTSCHER RS 300}
	{Nur mit angeschlossener Absaugung schleifen.}

\festool{Exzenterschleifer}{Festool ROTEX RO 150}
	{Nur mit angeschlossener Absaugung schleifen.}

\festool[M004]{Oszillierer}{Festool VECTURO OS 400}
	{Werkstück fest einspannen.}

\festool{Anschlagschienen\-zubehör}{Festool}
	{Zubehör für die Führungs- und Anschlagschienen.}

\begin{schild}[farbe=gruen, modell=für ROTEX und RUTSCHER, einweisung=werkstatt-einweisung]
	{Schleifpapier}{Rund, viereckig und dreieckig}
Verbrauchtes Schleifpapier bitte entsorgen.

\textbf{Fehlt was?} Fast ausverkauft? Schreib es auf die Nachkaufliste am Whiteboard!
\end{schild}

\begin{schild}[farbe=gruen, modell=Dremel 3000, seite=dremel]
	{Multifunktionswerkzeug}{Nur für feine Arbeiten}
Für Gröberes haben wir den Fein-Bohrschleifer.

Wende dich bei Unklarheiten bitte an die Betreuer:innen!
\end{schild}

\begin{schild}[farbe=gruen, symbole={W017}]
	{Heißluftfön}{Benutzung mit Verstand}
Unbedingt nach Benutzung auf \texttt{Stufe 1} laufend abkühlen lassen.

Wende dich bei Unklarheiten an die Betreuer:innen.
\end{schild}

\begin{schild}[farbe=gelb, einweisung=ultraschallbad-einweisung,
	kontakt={FabLab-Aktive, \mail{technik@fablab.fau.de}}]
	{Ultraschallbad}{}
Bitte wende dich an die Betreuer:innen.
\end{schild}

% ---------------------------------------------------------------------
% Textil
% ---------------------------------------------------------------------
\begin{schild}[farbe=gelb, modell=Brother VR, einweisung=stickmaschine-einweisung]
	{Stickmaschine}{}
Bei Interesse frage die Betreuer:innen!
\end{schild}

\begin{schild}[farbe=gelb, modell=Pfaff 260, einweisung=naehmaschine-einweisung]
	{Nähmaschine}{}
Bei Interesse frage die Betreuer:innen!
\end{schild}

\begin{schild}[farbe=gruen, seite=oesen-und-druckknopfpresse]
	{Druckknopf- und Ösenpresse}{Bitte nicht mit Gewalt auf den Hebel drücken}
Niemals einen Hammer oder Ähnliches verwenden!

Wende dich bei Unklarheiten bitte an die Betreuer:innen!
\end{schild}

% ---------------------------------------------------------------------
% Material, Lager, Sonstiges
% ---------------------------------------------------------------------
\begin{schild}[farbe=gruen, seite=kosten]
	{Elektro- und Mechanikmaterial}{Selbstbedienung}
Fehlt was? Fast ausverkauft? Schreib es auf die Nachkaufliste am Whiteboard!

Die Preise stehen auf den Tütchen/Schubladen. Wenn eine Schublade ganz ausverkauft ist,
stelle sie bitte auf den Kopf.

Manche Mechanik-Teile, z.\,B. Muttern, sind in größerer Menge in der Nachschubkiste vorrätig.
\end{schild}

\begin{schild}[farbe=gruen, modell=FabLab-Kamera]
	{Kamera}{\werkstattEinweisung}
Gerne für Fotos von Projekten verwenden!

Fotos nach dem Übertragen von der Speicherkarte löschen.
Akku nach Benutzung aufladen und die Kamera vollständig zurücklegen.
\end{schild}

\begin{schild}[farbe=gruen]
	{Fundsachenkiste}{Alles mit \enquote{FFA} Beschriftete ist zu verschenken. Der Rest sind Fundsachen.}
Hier landet der auf Tischen liegengelassene Schrott. Die Kiste wird nach jedem
FabLab-Treffen stark aussortiert und der Großteil weggeschmissen.
\end{schild}

% Generisches Schild
\begin{schild}[farbe=rot, modell={Gerät: \rule{6cm}{0.4pt}}, startseite]
	{Außer Betrieb}{Darf nicht benutzt werden}
Es wurde noch keine Einweisung erstellt und freigegeben.

\textbf{Bei Zuwiderhandlung erfolgt Lab-Verbot!}
\end{schild}
